\section*{Chapter \ref{ch:g10:tuples-sets-comprehensions} --- Tuples, Sets and Comprehensions}

\begin{solution}{exo:g10:tuples-sets-comprehensions:1}
After the first line \code{a}, \code{b}, \code{c} are $1$, $2$, $3$; the second line builds the \omterm{def:g10:tuples-sets-comprehensions:tuple}{tuple} \code{(3, 1)}
and gives $3$ to \code{a} and $1$ to \code{c}: the line prints \code{3 2 1}. With \code{x, y = (1, 2, 3)} there are
three values for two names: Python stops with \code{ValueError: too many values to unpack}.
\end{solution}

\begin{solution}{exo:g10:tuples-sets-comprehensions:2}
\code{s & t} is \code{{3, 4}}, \code{s | t} is \code{{1, 2, 3, 4, 5}}, \code{s - t} is \code{{1, 2}},
\code{t - s} is \code{{5}}, \code{len(s | t)} is $5$, and \code{5 in s} is \code{False}.
\end{solution}

\begin{solution}{exo:g10:tuples-sets-comprehensions:3}
\code{def first_last(t):} with the \omterm{def:g10:specifications-and-tests:docstring}{docstring} ``Return the first and the last element of t. \omterm{def:g10:specifications-and-tests:pre-post}{Precondition}: t is a
non-empty list.'' and the body \code{return t[0], t[-1]}. A call: \code{first, last = first_last([4, 8, 1])} gives
\code{first} the value $4$ and \code{last} the value $1$. For a list of one element both are that element.
\end{solution}

\begin{solution}{exo:g10:tuples-sets-comprehensions:4}
\code{u} is an \omterm{def:g10:tuples-sets-comprehensions:alias}{alias} of \code{t}; \code{v} is a copy. The line prints
\code{[1, 2, 3] [1, 2, 3] [1, 2, 4] True False True}. The drawing has two lists: \code{t} and \code{u} point to
\code{[1, 2, 3]}, \code{v} to \code{[1, 2, 4]}.
\end{solution}

\begin{solution}{exo:g10:tuples-sets-comprehensions:5}
\code{len(a) + len(b)} counts every element of the union once, except those of the intersection, which belong to
both \omterm{def:g10:tuples-sets-comprehensions:set}{sets} and are counted twice; subtracting \code{len(a & b)} counts them once. On the town run: $6 + 11 - 5 = 12$,
the size printed for the union. The union has \code{len(a) + len(b)} elements exactly when \code{a & b} is empty.
\end{solution}

\begin{solution}{exo:g10:tuples-sets-comprehensions:6}
\code{("red")} is only the string \code{"red"} in parentheses, and a loop over a string goes through its characters:
it prints \code{r}, \code{e}, \code{d} on three lines. The \omterm{def:g10:tuples-sets-comprehensions:tuple}{tuple} of one element is written \code{("red",)}, with the
comma.
\end{solution}

\begin{solution}{exo:g10:tuples-sets-comprehensions:7}
The file \code{exercises/common_books.py} builds the \omterm{def:g10:tuples-sets-comprehensions:set}{set} of the books of member \code{a} as

\smallskip\noindent\hspace*{2em}\code{books_a = {loan["book_id"] for loan in loans if loan["member_id"] == a}}

\smallskip\noindent and \code{books_b} the same way with \code{b}; the function returns \code{books_a & books_b}. For members $40$ and $55$ it returns \code{{16, 57, 62, 71, 75}}, as the
town listing printed (checked by the chapter's tests).
\end{solution}

\begin{solution}{exo:g10:tuples-sets-comprehensions:8}
\code{[0] * 3} builds one list \code{[0, 0, 0]}; multiplying the outer list by $3$ repeats the \emph{reference} to
it, so the grid is an outer list of three arrows to the same row. Writing through \code{grid[0]} changes that row,
which is also \code{grid[1]} and \code{grid[2]}. The \omterm{def:g10:tuples-sets-comprehensions:comprehension}{comprehension} \code{[[0] * 3 for _ in range(3)]} evaluates
\code{[0] * 3} three times and builds three rows (\code{exercises/grid_bug.py}).
\end{solution}

\begin{solution}{exo:g10:tuples-sets-comprehensions:9}
\code{{"a": "A", "b": "B"}}; then the \omterm{def:g10:tuples-sets-comprehensions:set}{set} of the remainders, \code{{0, 1, 2}}; then
\code{{0: 0, 1: 1, 3: 9}}, where the key $2$ is left out by the condition.
\end{solution}

\begin{solution}{exo:g10:tuples-sets-comprehensions:10}
A subset is fixed by deciding, for each of the four elements, whether it is in or out: $2$ choices for \code{"a"},
times $2$ for \code{"b"}, and so on, so $2 \times 2 \times 2 \times 2 = 16$ subsets, from the empty \omterm{def:g10:tuples-sets-comprehensions:set}{set} to the whole
\omterm{def:g10:tuples-sets-comprehensions:set}{set}. The subsets of two elements are the pairs \code{ab}, \code{ac}, \code{ad}, \code{bc}, \code{bd}, \code{cd}: $6$.
The \omterm{def:g10:tuples-sets-comprehensions:comprehension}{comprehension} of \code{exercises/subsets.py} lists the $16$ (checked by the chapter's tests).
\end{solution}

\begin{solution}{exo:g10:tuples-sets-comprehensions:11}
The function \code{distinct_authors} is one \omterm{def:g10:tuples-sets-comprehensions:set}{set} \omterm{def:g10:tuples-sets-comprehensions:comprehension}{comprehension},

\smallskip\noindent\hspace*{2em}\code{return {book["author"] for book in books}}

\smallskip\noindent and \code{authors_in_genre} adds the condition \code{if book["genre"] == genre} to it and
passes the \omterm{def:g10:tuples-sets-comprehensions:set}{set} to \code{sorted} (\code{exercises/distinct_authors.py}). The town library's $80$ books have $23$ distinct authors.
\end{solution}

\begin{solution}{exo:g10:tuples-sets-comprehensions:12}
\emph{(a)} Each of the $6$ elements of \code{a} is compared with at most the $11$ elements of \code{b}: at most
$6 \times 11 = 66$ comparisons. \emph{(b)} It is reached when no element of \code{a} is found before the end of
\code{b}: for example when the lists have nothing in common. \emph{(c)} $n \times n = n^2$; when $n$ doubles, the
largest number is multiplied by $4$. With \omterm{def:g10:tuples-sets-comprehensions:set}{sets}, \code{x in s} does not go through \code{s}, which is why the \omterm{def:g10:tuples-sets-comprehensions:set}{set}
version is much faster on large lists (\code{exercises/count_intersection.py} counts the comparisons).
\end{solution}

\begin{solution}{exo:g10:tuples-sets-comprehensions:13}
\code{return [[row[j] for row in table] for j in range(len(table[0]))]}: the outer \omterm{def:g10:tuples-sets-comprehensions:comprehension}{comprehension} goes through the
column numbers, the inner one collects the value of that column in each row. On \code{[[1, 2, 3], [4, 5, 6]]} it gives
\code{[[1, 4], [2, 5], [3, 6]]}, and transposing that gives the table back (\code{exercises/transpose.py}).
\end{solution}

\begin{solution}{exo:g10:tuples-sets-comprehensions:14}
\code{list(grid)} is a \omterm{def:g10:tuples-sets-comprehensions:copy}{shallow copy}: \code{new} is a new outer list, but its rows are the rows of \code{grid}. Writing
\code{new[r][c]} writes into the caller's rows, so the ``undo'' state changes with the new one. The fix is a \omterm{def:g10:tuples-sets-comprehensions:copy}{deep copy},
\code{new = copy.deepcopy(grid)}, or the \omterm{def:g10:tuples-sets-comprehensions:comprehension}{comprehension} \code{new = [list(row) for row in grid]}, which copies each
row (\code{exercises/grid_copy.py}; the tests check that the caller's grid is untouched).
\end{solution}

\begin{solution}{exo:g10:tuples-sets-comprehensions:15}
In \code{exercises/jaccard.py}: \code{assert a | b, "two empty sets have no Jaccard index"}, then
\code{return Fraction(len(a & b), len(a | b))}. The intersection is part of the union, so the index lies between
$0$ and $1$. It is $0$ when the \omterm{def:g10:tuples-sets-comprehensions:set}{sets} have no element in common, and $1$ when the intersection is the whole union, that
is, when the two \omterm{def:g10:tuples-sets-comprehensions:set}{sets} are equal.
\end{solution}

\begin{solution}{pb:g10:tuples-sets-comprehensions:1}
\textbf{1.} \code{books.setdefault(m, set())} returns \code{books[m]} if the key \code{m} is already there; otherwise
it first stores an empty \omterm{def:g10:tuples-sets-comprehensions:set}{set} under \code{m} and returns it. Then \code{.add(b)} puts the book in that \omterm{def:g10:tuples-sets-comprehensions:set}{set}. A \omterm{def:g10:tuples-sets-comprehensions:set}{set} is
used because a book borrowed twice must count once, and because \code{&} and \code{|} are needed next.

\textbf{2.} $60$ members borrowed at least one book, and $80$ different books were borrowed.

\textbf{3.} A member who borrowed the same book twice has it once in her \omterm{def:g10:tuples-sets-comprehensions:set}{set}. On the town data the sizes add up to
$388$: twelve loans repeat a book already borrowed by the same member.

\textbf{4.} With \code{ids = {l["member_id"] for l in loans}}, the dictionary is
\code{{m: {l["book_id"] for l in loans if l["member_id"] == m} for m in ids}}. For each of the $60$ members, the inner \omterm{def:g10:tuples-sets-comprehensions:comprehension}{comprehension} goes through all $400$ loans: $24{,}000$ steps
instead of the $400$ of the loop.

\textbf{5.} Sizes $6$ and $11$; intersection $5$; union $12$.

\textbf{6.} $5/12 \approx 0.417$.

\textbf{7.} With herself: $1$, since a \omterm{def:g10:tuples-sets-comprehensions:set}{set} and itself have the same intersection and union. With no common book: $0$.

\textbf{8.} Intersection $1$, union $2$: the index is $1/2$.

\textbf{9.} Without it, each pair would appear twice (\code{i, j} and \code{j, i}), and each member would be paired
with herself, with the index $1$.

\textbf{10.} Each of the $n$ members is paired with the $n - 1$ others, and each pair is then counted twice:
$n(n-1)/2$ pairs, that is $60 \times 59 / 2 = 1{,}770$.

\textbf{11.} \omterm{def:g10:tuples-sets-comprehensions:tuple}{Tuples} compare position by position, like words in a dictionary. The first component is the index with
a minus sign, so a larger index comes first; when two pairs have the same index, the member ids decide, so the order
never depends on chance.

\textbf{12.} Members $32$ and $47$, with $1/2 = 0.500$.

\textbf{13.} It rests on one shared book: member $32$ borrowed a single book. With so few books, one coincidence
gives a large index. Indices of very small \omterm{def:g10:tuples-sets-comprehensions:set}{sets} are not evidence of a common taste.

\textbf{14.} $49$ members and $49 \times 48 / 2 = 1{,}176$ pairs.

\textbf{15.} Members $40$ and $55$, with $5/12 \approx 0.417$.

\textbf{16.} Members $15$ and $56$, and members $31$ and $43$, both with $1/3 \approx 0.333$, then $41$ and $50$ with
$2/7 \approx 0.286$.

\textbf{17.} The intersection has at most $6$ elements (it is part of the smaller \omterm{def:g10:tuples-sets-comprehensions:set}{set}), and the union at least $11$
(it contains the larger one): the index is at most $6/11 \approx 0.545$. It reaches it exactly when the \omterm{def:g10:tuples-sets-comprehensions:set}{set} of $6$ is
contained in the \omterm{def:g10:tuples-sets-comprehensions:set}{set} of $11$. The pair $40$ and $55$ misses it by one book: one of Dara's six books is not among
Umar's.

\textbf{18.} One per pair: $1{,}176$.

\textbf{19.} The maximum over \code{itertools.combinations(ids, 2)} is again $5/12$ (checked by the chapter's tests).

\textbf{20.} Among the members with at least $5$ books, the two who read most alike are Dara Reedfield (member $40$,
$6$ books) and Umar Daleton (member $55$, $11$ books): they share $5$ books, and their Jaccard index is
$5/12 \approx 0.417$, ahead of the next pairs at $1/3$.
\end{solution}
